% GENERATED FILE. Edit canonical lessons, then run npm run generate:books.
% canonical-source-hash: fnv1a64:9a60a7702c165718
% canonical-lessons: MR-C80-visarne, MR-C80-pathavne, MR-C80-mojne, MR-C80-pakadne, MR-C80-kapne

\chapter{To forget, To send, To count, To catch, To cut}
\label{ch:mr-to-forget-to-send-to-count-to-catch-to-cut}
\hlchaptermodality{80}

\begin{chapteropening}
\textbf{By the end of this chapter:} \emph{I can say to forget, to send, to count, to catch and to cut.}

Complete the last lesson of chapter 80: I can say to forget, to send, to count, to catch and to cut.
\end{chapteropening}

\section[\texorpdfstring{\mr{विसरणे} (visarṇe)}{विसरणे (visarṇe)}]{\mr{विसरणे} (visarṇe) — to forget}

\label{lesson:MR-C80-visarne}



\begin{quote}
Before the new one: say the Marathi for to sing; a song, then the Marathi for to dance.
\end{quote}

\subsection*{You'll want to know: \mr{विसरणे}}
\textbf{\mr{विसरणे}} — \emph{visarṇe} — \textquotedblleft{}to forget\textquotedblright{}.

\begin{cousinweb}
\textbf{\mr{विसरणे}} (\emph{visarṇe}) is a tadbhava, a word inherited from Sanskrit and reshaped by sound change. It comes from \emph{vismarati}, \textquotedblleft{}forgets\textquotedblright{} (CDIAL 12021). It is \emph{vi-}, \textquotedblleft{}away\textquotedblright{}, and \emph{smṛ}, \textquotedblleft{}to remember\textquotedblright{}.
\end{cousinweb}

\begin{grammarlens}[title={What you've built}]
One more word for this chapter.
\end{grammarlens}

\subsection*{Your turn}
\begin{itemize}
\raggedright
  \item \emph{Say it:} \emph{visarṇe}
  \item \emph{Say it:} \emph{visarṇe}, once more
  \item \emph{Say it:} say \emph{visarṇe}
  \item \emph{Recall:} say the Marathi for to sing; a song, then the Marathi for to dance, then say \emph{visarṇe} again
\end{itemize}

\begin{tcolorbox}[breakable,colback=teal!4,colframe=teal!35!black,title={Before you move on}]
What does \mr{विसरणे} mean? (\textquotedblleft{}To forget\textquotedblright{}.) Say it once more.
\end{tcolorbox}
\section[\texorpdfstring{\mr{पाठवणे} (pāṭhavṇe)}{पाठवणे (pāṭhavṇe)}]{\mr{पाठवणे} (pāṭhavṇe) — to send}

\label{lesson:MR-C80-pathavne}



\begin{quote}
Before the new one: say the Marathi for to dance, then the Marathi for to forget.
\end{quote}

\subsection*{You'll want to know: \mr{पाठवणे}}
\textbf{\mr{पाठवणे}} — \emph{pāṭhavṇe} — \textquotedblleft{}to send\textquotedblright{}.

\begin{cousinweb}
\textbf{\mr{पाठवणे}} (\emph{pāṭhavṇe}) is a tadbhava, a word inherited from Sanskrit and reshaped by sound change. It comes from \emph{prasthāpayati}, \textquotedblleft{}puts aside, sends\textquotedblright{}, built on \emph{sthā}, \textquotedblleft{}to stand\textquotedblright{} (CDIAL 8607).
\end{cousinweb}

\begin{grammarlens}[title={What you've built}]
One more word for this chapter.
\end{grammarlens}

\subsection*{Your turn}
\begin{itemize}
\raggedright
  \item \emph{Say it:} \emph{pāṭhavṇe}
  \item \emph{Say it:} \emph{pāṭhavṇe}, once more
  \item \emph{Say it:} say \emph{pāṭhavṇe}
  \item \emph{Recall:} say the Marathi for to dance, then the Marathi for to forget, then say \emph{pāṭhavṇe} again
\end{itemize}

\begin{tcolorbox}[breakable,colback=teal!4,colframe=teal!35!black,title={Before you move on}]
What does \mr{पाठवणे} mean? (\textquotedblleft{}To send\textquotedblright{}.) Say it once more.
\end{tcolorbox}
\section[\texorpdfstring{\mr{मोजणे} (mojṇe)}{मोजणे (mojṇe)}]{\mr{मोजणे} (mojṇe) — to count}

\label{lesson:MR-C80-mojne}



\begin{quote}
Before the new one: say the Marathi for to forget, then the Marathi for to send.
\end{quote}

\subsection*{You'll want to know: \mr{मोजणे}}
\textbf{\mr{मोजणे}} — \emph{mojṇe} — \textquotedblleft{}to count\textquotedblright{}.

\begin{cousinweb}
\textbf{\mr{मोजणे}} (\emph{mojṇe}) is listed by CDIAL under Sanskrit \emph{meya}, \textquotedblleft{}to be measured\textquotedblright{}, from \emph{mā}, \textquotedblleft{}to measure\textquotedblright{}, beside Marathi \emph{mejṇẽ}, \textquotedblleft{}to measure\textquotedblright{}. CDIAL leaves the \emph{-o-} of \textbf{\mr{मोजणे}} unexplained (CDIAL 10329).
\end{cousinweb}

\begin{grammarlens}[title={What you've built}]
One more word for this chapter.
\end{grammarlens}

\subsection*{Your turn}
\begin{itemize}
\raggedright
  \item \emph{Say it:} \emph{mojṇe}
  \item \emph{Say it:} \emph{mojṇe}, once more
  \item \emph{Say it:} say \emph{mojṇe}
  \item \emph{Recall:} say the Marathi for to forget, then the Marathi for to send, then say \emph{mojṇe} again
\end{itemize}

\begin{tcolorbox}[breakable,colback=teal!4,colframe=teal!35!black,title={Before you move on}]
What does \mr{मोजणे} mean? (\textquotedblleft{}To count\textquotedblright{}.) Say it once more.
\end{tcolorbox}
\section[\texorpdfstring{\mr{पकडणे} (pakaḍṇe)}{पकडणे (pakaḍṇe)}]{\mr{पकडणे} (pakaḍṇe) — to catch}

\label{lesson:MR-C80-pakadne}



\begin{quote}
Before the new one: say the Marathi for to send, then the Marathi for to count.
\end{quote}

\subsection*{You'll want to know: \mr{पकडणे}}
\textbf{\mr{पकडणे}} — \emph{pakaḍṇe} — \textquotedblleft{}to catch\textquotedblright{}.

\begin{cousinweb}
\textbf{\mr{पकडणे}} (\emph{pakaḍṇe}) comes from a reconstructed \emph{*pakkaḍ-}, \textquotedblleft{}seize\textquotedblright{} (CDIAL 7619). CDIAL calls its link with Sanskrit \emph{pragṛhṇāti}, \textquotedblleft{}holds\textquotedblright{}, obscure and doubtful.
\end{cousinweb}

\begin{grammarlens}[title={What you've built}]
One more word for this chapter.
\end{grammarlens}

\subsection*{Your turn}
\begin{itemize}
\raggedright
  \item \emph{Say it:} \emph{pakaḍṇe}
  \item \emph{Say it:} \emph{pakaḍṇe}, once more
  \item \emph{Say it:} say \emph{pakaḍṇe}
  \item \emph{Recall:} say the Marathi for to send, then the Marathi for to count, then say \emph{pakaḍṇe} again
\end{itemize}

\begin{tcolorbox}[breakable,colback=teal!4,colframe=teal!35!black,title={Before you move on}]
What does \mr{पकडणे} mean? (\textquotedblleft{}To catch\textquotedblright{}.) Say it once more.
\end{tcolorbox}
\section[\texorpdfstring{\mr{कापणे} (kāpṇe)}{कापणे (kāpṇe)}]{\mr{कापणे} (kāpṇe) — to cut}

\label{lesson:MR-C80-kapne}



\begin{quote}
Before the new one: say the Marathi for to count, then the Marathi for to catch.
\end{quote}

\subsection*{You'll want to know: \mr{कापणे}}
\textbf{\mr{कापणे}} — \emph{kāpṇe} — \textquotedblleft{}to cut\textquotedblright{}.

\begin{cousinweb}
\textbf{\mr{कापणे}} (\emph{kāpṇe}) is a tadbhava, a word inherited from Sanskrit and reshaped by sound change. It comes from \emph{kalpayati}, \textquotedblleft{}sets in order; trims, cuts\textquotedblright{} (CDIAL 2944).
\end{cousinweb}

\begin{grammarlens}[title={What you've built}]
One more word for this chapter.
\end{grammarlens}

\subsection*{Your turn}
\begin{itemize}
\raggedright
  \item \emph{Say it:} \emph{kāpṇe}
  \item \emph{Say it:} \emph{kāpṇe}, once more
  \item \emph{Say it:} say \emph{kāpṇe}
  \item \emph{Recall:} say the Marathi for to count, then the Marathi for to catch, then say \emph{kāpṇe} again
\end{itemize}

\begin{tcolorbox}[breakable,colback=teal!4,colframe=teal!35!black,title={Before you move on}]
What does \mr{कापणे} mean? (\textquotedblleft{}To cut\textquotedblright{}.) Say it once more.
\end{tcolorbox}
